
\documentclass{article}
\usepackage{colortbl}
\usepackage{makecell}
\usepackage{multirow}
\usepackage{supertabular}

\begin{document}

\newcounter{utterance}

\centering \large Interaction Transcript for game `ifeval', experiment `ifeval', episode 62 with llama{-}3.1{-}8b{-}instruct.
\vspace{24pt}

{ \footnotesize  \setcounter{utterance}{1}
\setlength{\tabcolsep}{0pt}
\begin{supertabular}{c@{$\;$}|p{.15\linewidth}@{}p{.15\linewidth}p{.15\linewidth}p{.15\linewidth}p{.15\linewidth}p{.15\linewidth}}
    \# & \multicolumn{2}{c}{Player} && \multicolumn{2}{c}{Game Master} \\
    \hline

    \theutterance \stepcounter{utterance}  
    & & & \multicolumn{4}{p{0.6\linewidth}}{
        \cellcolor[rgb]{0.9,0.9,0.9}{
            \makecell[{{p{\linewidth}}}]{
                \texttt{\tiny{[P1$\langle$GM]}}
                \texttt{Which of the following is not a fish: salmon or avocado? Answer this easy question first, then expand into an interesting article about salmon and avocado. Use some words in all caps in your response, but those all{-}caps words should appear at most 10 times.} \\
            }
        }
    }
     \\ \\

    \theutterance \stepcounter{utterance}  
    & \multicolumn{4}{p{0.6\linewidth}}{
        \cellcolor[rgb]{0.561,0.737,0.561}{
            \makecell[{{p{\linewidth}}}]{
                \texttt{\tiny{[P1$\rangle$GM]}}
                \texttt{The answer is clearly AVOCADO. Avocado is a type of fruit, not a fish.} \\
\\ 
\texttt{Now, let's dive into an interesting article about SALMON and AVOCADO.} \\
\\ 
\texttt{**The Amazing World of SALMON and AVOCADO**} \\
\\ 
\texttt{SALMON, a type of ANADROMOUS fish, has been a staple in many cuisines around the world for centuries. These fish are known for their rich, FATTY ACID content, which makes them a popular choice for health{-}conscious individuals. SALMON is also an excellent source of PROTEIN, VITAMIN D, and OMEGA{-}3 FATTY ACIDS, making it a great addition to a balanced diet.} \\
\\ 
\texttt{On the other hand, AVOCADO is a nutrient{-}dense fruit that has gained immense popularity in recent years. This GREEN, CREAMY fruit is rich in HEALTHY FATS, FIBER, and various essential VITAMINS and MINERALS. AVOCADO is also a great source of ANTIOXIDANTS, which help protect the body against FREE RADICALS and oxidative stress.} \\
\\ 
\texttt{Interestingly, both SALMON and AVOCADO have been linked to several health benefits. SALMON has been shown to reduce the risk of HEART DISEASE, while AVOCADO has been linked to improved BLOOD SUGAR CONTROL and WEIGHT MANAGEMENT. Additionally, both SALMON and AVOCADO are rich in ANTINFLAMMATORY compounds, which may help reduce inflammation and improve overall health.} \\
\\ 
\texttt{In terms of culinary uses, SALMON and AVOCADO are both versatile ingredients that can be used in a variety of dishes. SALMON can be grilled, baked, or poached, while AVOCADO can be used in salads, smoothies, or as a topping for various dishes. When combined, SALMON and AVOCADO make for a delicious and nutritious pairing, perfect for health{-}conscious individuals looking to add some excitement to their meals.} \\
\\ 
\texttt{In conclusion, while SALMON and AVOCADO may seem like an unlikely pair, they share many similarities in terms of their nutritional benefits and culinary uses. Whether you're a fan of FISH or FRUIT, both SALMON and AVOCADO are definitely worth incorporating into your diet.} \\
            }
        }
    }
    & & \\ \\

    \theutterance \stepcounter{utterance}  
    & & & \multicolumn{2}{p{0.3\linewidth}}{
        \cellcolor[rgb]{0.9,0.9,0.9}{
            \makecell[{{p{\linewidth}}}]{
                \texttt{\tiny{[GM$|$GM]}}
                {'change_case:capital_word_frequency': {'capital_relation': 'at least', 'capital_frequency': 1}}
            }
        }
    }
    & & \\ \\

    \theutterance \stepcounter{utterance}  
    & & & \multicolumn{2}{p{0.3\linewidth}}{
        \cellcolor[rgb]{0.9,0.9,0.9}{
            \makecell[{{p{\linewidth}}}]{
                \texttt{\tiny{[GM$|$GM]}}
                \texttt{game\_result = WIN} \\
            }
        }
    }
    & & \\ \\

\end{supertabular}
}

\end{document}
